\pagestyle{fancy} 
\section{Future Work}\label{sec:future-work}

The present evaluation is intentionally narrow: one scenario, one operationalisation of the framework, and a remote, self-paced procedure. The directions below follow from that scope and from what the artifact already suggests about where ephemeral, task-scoped support may be most informative next.

\subsection{Broader scenarios and user populations}

The study used a single presentation-refinement workflow. Ephemeral support is unlikely to matter equally in every domain. Prior work like with BISCUIT \parencite{Cheng2024Biscuit} suggest that ephemeral interface scaffolds work well with computational notebooks for data science exploration where they can steer intermediate output. Thus, it provides a useful contrast to this thesis where it focused on specific technical workflow, whereas this projects focuses if ephemeral support can add value in conventional interfaces for a larger user demographic. The current results show that less focus should be made on broad domains and more on the workflow properties of interruption tolerance, consequence of error, and whether the assisted action is a bounded local refinement rather than a high-stakes commitment. The most interesting scenarios to test next would be technical refinement tasks through several iterations. For example, when a designer wants to tweak the layout or hierarchy of their designs or a data scientist adjusting the parameters of a machine learning model they both require real time feedback from their actions. The same reasoning indicates significant limitations regarding the appropriateness of ephemeral support. Within heavily regulated areas like finance, law, or healthcare ephemeral generative UI may not be an ideal use case as those scenarios. This is as they require stricter predictability, traceability, and accountability than low risk situations like adding tasks to a todo list, even if it's proven useful for refinement tasks. This possibility also aligns with the claim that adaptive support could be the most valuable when traditional interfaces struggle to accommodate varying user needs across different contexts \parencite{Bieniek2024GenerativeAI,Tambi2020GenerativeBanking}. The next step is to evaluate the framework across several workflows that differ in interruption tolerance, consequence of error, task significance, technical complexity, familiarity with the application, and degree of regulatory constraint. It would clarify situations when ephemeral support is genuinely helpful for end users or if it gives irrelevant guidance for a clear user workflow.

\subsection{Deeper evaluation methods}


Increasing the number of completions would make the statistics more reliable, but it wouldn't explain why participants acted the way they did or how they understood the support. Another direction would be to do in depth evaluation in focus groups where a researcher can observe the misunderstanding of triggers and emotional reactions of the participants that questionnaires can only partially capture. Using both controlled experiments and focus groups could link specific design choices to real friction and also to how intrusive or controllable the support felt.

\subsection{Richer ephemeral component catalogues}

The evaluation artifact deliberately used a constrained set of ephemeral UI components so that conditions remained comparable and the study remained feasible within the capstone timeline. The underlying framework could support a wider catalogue: more varied layouts, stronger personalisation of content and tone based on the user's profile, role, or prior behaviour, and more expressive temporary interaction patterns than were deployed here. Exploring that wider design space would also connect more directly to the broader literature, which frames generative interfaces in terms of adaptive, context-sensitive interaction structures while also warning that flexibility can come at the cost of predictability and user control \parencite{Bieniek2024GenerativeAI}. At the same time, the foundational logic of ephemerality suggests that temporary interface elements should appear only when relevant and should recede without leaving persistent clutter \parencite{Doring2013EphemeralUI}. Expanding the catalogue therefore raises its own research questions about how to preserve predictability and user agency when surface behaviour varies more widely, and how to evaluate more personalised variants without confounding condition effects with implementation novelty. Future work could treat catalogue design as a first-class object of study: progressive disclosure of complexity, user-tunable boundaries for ephemerality, and systematic comparison of a small number of high-diversity implementations against a stable baseline.
